\subsection{UAV Telemetry Anomaly Detection and Fault Diagnosis}

UAV anomaly detection and fault diagnosis aim to identify abnormal operating states from flight measurements and determine the corresponding fault conditions. Existing approaches can be broadly categorized as knowledge-based, model-based, signal-based, and data-driven methods \cite{fourlas2021survey,puchalski2022uav,adaika2025fdd}. With the increasing availability of onboard flight logs, data-driven approaches have received growing attention because multivariate telemetry records the coupled behavior of state estimation, inertial sensing, actuation, power systems, and other flight-control components. Real-flight datasets such as ALFA and UAV-SEAD further provide annotated telemetry for evaluating anomaly detection and fault diagnosis under realistic operating conditions \cite{keipour2021alfa,kabaoglu2026uavsead,keipour2019automatic}.

Early data-driven studies mainly relied on conventional machine-learning techniques to identify deviations from normal flight behavior. More recent methods increasingly exploit temporal and cross-channel dependencies using reconstruction models, probabilistic sequence models, graph neural networks, and attention-based architectures \cite{zhang2019mscred,su2019omnianomaly,deng2021gdn,tuli2022tranad,xu2022anomaly}. In UAV-specific anomaly detection, MTCL-UAV employs adaptive multiscale Transformer modeling and contrastive learning to capture temporal variations and correlations in flight data \cite{hu2025mtcl}, while LDC-P-VAE introduces dynamics-constrained probabilistic modeling for adaptive anomaly detection \cite{yang2026ldcpvae}. These methods demonstrate the effectiveness of learning richer temporal representations from UAV telemetry and have substantially advanced flight-data anomaly detection.

Despite this progress, most existing methods primarily produce anomaly scores, binary alarms, or predefined fault labels. Such outputs are effective for determining whether abnormal behavior occurs, but provide limited support for subsequent engineering inspection. In particular, the model evidence behind a fault prediction is rarely organized into physical telemetry channels or connected further to the flight-control software architecture. Our work builds on UAV telemetry anomaly detection and fault diagnosis by extending the diagnostic process from abnormal-state identification toward interpretable telemetry evidence and architecture-aware software inspection \cite{lundberg2017shap,lundberg2020trees,meier2015px4,px4uorbdocs2026}.
